El moviment dels insectes en la teranyina feta per les aranyes és un moviment harmònic simple (MHS), és a dir, es pot modelitzar com una massa a l'extrem d'una molla. S'ha observat que quan l'aranya està sola a la teranyina produeix una vibració de freqüència $12\ \mathrm{Hz}$. Si un insecte d'$1,00\ \mathrm{g}$ de massa queda atrapat a la teranyina, el conjunt aranya i insecte produeix una vibració de $10\ \mathrm{Hz}$.

\begin{apartat}{1,25}
Calculeu la massa de l'aranya.
\end{apartat}

\begin{apartat}{1,25}
Calculeu la constant elàstica d'aquesta teranyina. En quines posicions aquest MHS assoleix la màxima velocitat? I la màxima acceleració?
\end{apartat}
